\documentclass[a4paper,10pt,landscape]{article}

\usepackage[czech]{babel}
\usepackage[T1]{fontenc}
\usepackage[utf8]{inputenc}
\usepackage{amsmath,amssymb}
\usepackage{geometry}
\usepackage{enumitem}

\geometry{a4paper,landscape,left=8mm,right=8mm,top=8mm,bottom=8mm}
\pagestyle{empty}
\setlength{\parindent}{0pt}
\setlength{\parskip}{0pt}
\setlist[enumerate]{leftmargin=*,itemsep=2pt,topsep=2pt}

% Uprav název testu; jednotlivé varianty jsou vždy A/B, C/D atd.
\newcommand{\TestTitle}{Geometrická posloupnost}
\newcommand{\TestHeader}[1]{%
    {\bfseries TEST -- \TestTitle}\hfill{\bfseries Varianta #1}\\[2mm]
    Jméno a příjmení: \rule{4.5cm}{0.4pt}\hfill Třída: \rule{1.5cm}{0.4pt}\\[2mm]
    Datum: \rule{2.4cm}{0.4pt}\hfill Počet bodů: \rule{1.4cm}{0.4pt}\hfill Známka: \rule{1.4cm}{0.4pt}\\[1mm]
    \hrule\vspace{3mm}
}
\newcommand{\AnswerSpace}[1]{\vspace{#1}}

\begin{document}
\noindent
\begin{minipage}[t]{0.48\textwidth}
    \small
    \TestHeader{A}

    % ZADÁNÍ VARIANTY A
    \begin{enumerate}
        \item Rozhodněte, zda čísla $5,15,45$ tvoří tři po sobě jdoucí členy geometrické posloupnosti. Pokud ano, určete kvocient. Odpověď zdůvodněte.
        \AnswerSpace{18mm}
        \item V geometrické posloupnosti platí $a_4=24$, $q=2$. Určete $a_1$ a $a_7$.
        \AnswerSpace{25mm}
        \item Určete součet prvních šesti členů geometrické posloupnosti $3,6,12,24,\ldots$.
        \AnswerSpace{25mm}
        \item Určete všechna $x\in\mathbb R$ tak, aby čísla $x-1$, $x+1$, $x+5$ tvořila v uvedeném pořadí tři po sobě jdoucí členy geometrické posloupnosti. Proveďte zkoušku.
        \AnswerSpace{25mm}
    \end{enumerate}
\end{minipage}%
\hfill\vrule width 0.4pt\hfill
\begin{minipage}[t]{0.48\textwidth}
    \small
    \TestHeader{B}

    % ZADÁNÍ VARIANTY B
    \begin{enumerate}
        \item Rozhodněte, zda čísla $7,21,63$ tvoří tři po sobě jdoucí členy geometrické posloupnosti. Pokud ano, určete kvocient. Odpověď zdůvodněte.
        \AnswerSpace{18mm}
        \item V geometrické posloupnosti platí $a_4=40$, $q=2$. Určete $a_1$ a $a_7$.
        \AnswerSpace{25mm}
        \item Určete součet prvních šesti členů geometrické posloupnosti $5,10,20,40,\ldots$.
        \AnswerSpace{25mm}
        \item Určete všechna $x\in\mathbb R$ tak, aby čísla $x-3$, $x$, $x+6$ tvořila v uvedeném pořadí tři po sobě jdoucí členy geometrické posloupnosti. Proveďte zkoušku.
        \AnswerSpace{25mm}
    \end{enumerate}
\end{minipage}
\end{document}
